% Part: computability
% Chapter: computability-theory
% Section: complete-ce-sets

\documentclass[../../../include/open-logic-section]{subfiles}

\begin{document}

\olfileid{cmp}{thy}{cce}
\olsection{সম্পূর্ণ গণনসাধ্যভাবে তালিকায়নযোগ্য সেট}

\begin{defn}
কোনো সেট~$A$ একটি \emph{সম্পূর্ণ !!{computably enumerable} সেট}
(বহু-এক হ্রাসযোগ্যতার অধীনে) যদি
\begin{enumerate}
\item $A$ গণনসাধ্যভাবে তালিকায়নযোগ্য হয়, এবং
\item অন্য যেকোনো গণনসাধ্যভাবে তালিকায়নযোগ্য সেট~$B$-এর জন্য
$B \leq_m A$।
\end{enumerate}
\end{defn}

অন্যভাবে বললে, সম্পূর্ণ গণনসাধ্যভাবে তালিকায়নযোগ্য সেটগুলি সম্ভাব্য
“কঠিনতম” গণনসাধ্যভাবে তালিকায়নযোগ্য সেট। এগুলি \emph{যেকোনো}
গণনসাধ্যভাবে তালিকায়নযোগ্য সেট-সম্পর্কিত প্রশ্নের উত্তর দিতে সাহায্য করে।

\begin{thm}
$K$, $K_0$ এবং $K_1$—সবগুলিই সম্পূর্ণ গণনসাধ্যভাবে তালিকায়নযোগ্য সেট।
\end{thm}

\begin{proof}
$K_0$ যে সম্পূর্ণ তা দেখতে $B$-কে যেকোনো গণনসাধ্যভাবে তালিকায়নযোগ্য
সেট ধরি। তাহলে কোনো সূচক~$e$-এর জন্য
\[
B = W_e = \Setabs{x}{\cfind{e}(x) \fdefined}.
\]
$f$-কে $f(x) = \tuple{e, x}$ অপেক্ষকটি ধরি। তাহলে প্রত্যেক স্বাভাবিক
সংখ্যা~$x$-এর জন্য $x \in B$ যদি এবং কেবল যদি $f(x) \in K_0$। অন্যভাবে
বললে, $f$, $B$-কে~$K_0$-তে হ্রাস করে।

$K_1$ যে সম্পূর্ণ তা দেখতে লক্ষ করো, \olref[k1]{prop:k1}-এর প্রমাণে
আমরা $K_0$-কে তার মধ্যে হ্রাস করেছি। তাই \olref[ppr]{prop:trans-red}
অনুসারে যেকোনো গণনসাধ্যভাবে তালিকায়নযোগ্য সেটকেও~$K_1$-তে হ্রাস করা যায়।

$K_0$-কে অনেকটা একইভাবে~$K$-তে হ্রাস করা যায়।
% BN-SRC-154 TARGET-CORRECTION-BEGIN
\par\smallskip\noindent\textnormal{\small\textbf{উৎস-সংশোধন (BN-SRC-154):} হিমায়িত ইংরেজি শেষ বাক্যে $K$-কে $K_0$-তে হ্রাস করার কথা বলা হয়েছে; সে দিকটি $K_0$-এর কঠিনতা সমর্থন করে, $K$-এর সম্পূর্ণতা নয়। এই অংশে প্রয়োজনীয় দিকটি $K_0 \leq_m K$; বাংলা পাঠে সেটিই রাখা হয়েছে।} % BN-SRC-154 TARGET-CORRECTION-END
\end{proof}

\begin{prob}
$K$ থেকে~$K_0$-তে একটি হ্রাসকরণ দাও।
\end{prob}

\begin{digress}
তাহলে দেখা গেল, এ পর্যন্ত বিবেচিত গণনসাধ্যভাবে তালিকায়নযোগ্য সেটের
সব উদাহরণ হয় গণনসাধ্য, নয়তো সম্পূর্ণ। এটি অদ্ভুত মনে হওয়া উচিত!{}
এমন কোনো গণনসাধ্যভাবে তালিকায়নযোগ্য সেট আছে কি, যা গণনসাধ্যও নয়,
সম্পূর্ণও নয়? উত্তর হ্যাঁ; কিন্তু 1950-এর দশকের মাঝামাঝি সময়ে
ফ্রিডবার্গ ও মুচনিক স্বাধীনভাবে এটি প্রতিষ্ঠা করার আগে তা জানা যায়নি।
\end{digress}

\end{document}
